\subsection{Local Affine Reduction and Event-Conditioned Controls}

For each focal individual and frame, candidate neighbors were first restricted
to individuals present at both the reference frame and the finite-lag frame.
These candidates were then ranked by Euclidean distance at the reference
frame. The focal individual was excluded, the nearest $k$ retained candidates
defined the local neighborhood, and focal samples with fewer than four
retained neighbors were discarded. The default lag was 0.10 s. The initial
survival analysis used two neighborhood sizes, $k=8$ and $k=10$, and later
focal-activity analyses used the frozen $k=8$ representation.

The local affine deformation was fitted by equal-weight least squares to
finite-lag relative neighbor displacements. Thus, although the raw trajectory
files include velocity fields, the frozen local T1 observable was computed
from relative displacement residuals divided by the lag duration. This choice
keeps the local affine baseline tied to the same finite-lag neighborhood used
to define the residual.

Local frame metrics were robust-z standardized within each observation. Slow
trends were removed by subtracting a one-second centered rolling mean, and the
resulting residual traces were robust-z standardized again before
event-conditioned comparisons. Because this detrending uses a centered window,
event-aligned time profiles should be interpreted as descriptive association
analyses rather than causal or online prediction analyses.

For the survival and spatial analyses, pre/post summaries used
$[-0.20,0)$ s and $[0,0.20]$ s windows relative to each transition time.
Same-observation non-event controls were generated by sampling one control
time per real event, requiring each control center to be at least 0.80 s away
from true transition times. Forty non-event replicates were used. A T1
survival call at a given neighborhood scale required all three frozen screens:
the local event-control excess
\texttt{local\_event\_minus\_non\_event\_direction\_z} had to exceed 0.03
robust-z units, the empirical non-event tail fraction
\texttt{p\_non\_event\_direction\_ge\_event} had to be no greater than 0.35,
and the local/global-affine retention ratio
\texttt{local\_to\_b3\_direction\_ratio} had to be at least 0.30.
Here, B3 denotes the upstream global-affine residual baseline, so this ratio
required the local residual to remain non-negligible relative to that
global-affine reference. The 0.35 tail criterion was a frozen screening
component rather than a conventional single-observation significance test. It
defined a consistent residual-support call from direction, empirical tail
rank, and retention relative to the global-affine residual. Statistical
calibration of the resulting all-observation both-scale count was performed
separately with the full-pipeline pseudo-event omnibus null.
The main all-observation count used a two-scale consistency requirement:
survival was counted as both-scale support only when the same observation
passed this per-scale screen at both original neighborhood sizes, $k=8$ and
$k=10$.

Three post-freeze checks were then used to harden the manuscript-facing
interpretation without changing the frozen T1 definition. First, a
full-pipeline pseudo-event omnibus null reran the same survival gate on
same-observation pseudo-event structures that avoided true transition windows.
The completed high-replicate run used $B=1000$ null replicates and 40
non-event controls per replicate-observation. For computational tractability,
the run was split into deterministic chunks and merged after completion; this
changed only execution, not the frozen pseudo-event design or survival gate.
Second, a detrending challenge reran the survival and timing summaries under
three fixed variants: robust-z standardization followed by centered one-second
rolling-mean subtraction and a second robust-z standardization; the same
procedure with a trailing, past-only one-second rolling mean; and a no-rolling
variant using robust-z standardization only. Third, local affine conditioning
quality control sampled frames across all observations and both original
neighborhood scales, recording valid-fit fractions, rank sufficiency,
condition numbers, and whether the largest T1 samples were concentrated in
ill-conditioned fits.

\subsection{Scale, Spatial, and Timing Analyses}

Surviving observations were rechecked under nearby choices of local scale and
lag. This analysis was used as a robustness check rather than as a search for
an optimal parameter set.

Spatial diagnostics decomposed unsigned T1 activity into diffuse,
radial-shell, and local-density summaries. Diffuse activity pooled the
tangential residual magnitude across sampled focal neighborhoods. Radial-shell
diagnostics divided focal individuals into core, middle, and edge thirds by
their swarm-centered radius. Local-density diagnostics used the distance to
the $k$-th retained neighbor as a density proxy and compared dense, middle,
and sparse thirds. Edge/core and dense/sparse contrasts were treated as
candidate spatial descriptions, not as predefined mechanisms.

Timing analyses used an event-aligned profile from $-0.50$ s to $0.50$ s
around each transition, sampled every 0.05 s. Four phase bins were used:
early-pre $[-0.50,-0.25)$ s, near-pre $[-0.25,0.00)$ s, near-post
$[0.00,0.25]$ s, and late-post $(0.25,0.50]$ s. Signed event-type analyses
then compared low-to-high and high-to-low transitions within these phase
windows. These signed analyses were treated as boundary tests because they
asked whether the unsigned T1 activity also had a stable directional or mirror
structure.
